\section{\benchmark: Budgeted Multi-Stream VLM Monitoring}

\benchmark{} converts existing event-centric video resources into online monitoring episodes. Each episode multiplexes many source videos into a synchronized stream set. Event intervals retain labels, severity, and natural-language descriptions. Non-event videos provide background streams and hard distractors.

\paragraph{Episode construction.}
Given a pool of annotated videos, we sample $N$ streams per episode, align them on a common clock, and place rare events at their original or resampled intervals. The agent observes only cheap per-stream features at every timestep. Expensive VLM outputs are hidden unless the policy spends a query. This creates an online decision problem rather than an offline video QA task.

\paragraph{Dataset tracks.}
The primary surveillance track uses UCF-Crime-style untrimmed anomaly videos. The robustness track adds XD-Violence-style multimodal violence events and smaller scene-specific anomaly datasets. The semantic track attaches open-vocabulary incident prompts and summary targets to selected event intervals. Long-video QA datasets can be adapted as a diagnostic track, but they are not treated as substitutes for rare-event monitoring.

\paragraph{Comparison to existing benchmarks.}
Table~\ref{tab:benchmark-comparison} summarizes the gap. Existing VQA, long-video, anomaly, and streaming-video benchmarks each cover part of the problem, but they do not jointly require multi-stream online allocation under explicit VLM budgets.

\begin{table}[t]
\centering
\caption{Benchmark families and whether they evaluate the properties required by budgeted multi-stream VLM monitoring.}
\label{tab:benchmark-comparison}
\begin{tabular}{lccccc}
\toprule
Benchmark family & Multi-stream & Online & Budgeted VLM & Event latency & Semantic reports \\
\midrule
Image/VQA benchmarks & \xmark & \xmark & \xmark & \xmark & \cmark \\
Long-video QA & \xmark & \xmark & \xmark & \xmark & \cmark \\
Video anomaly detection & \xmark & \cmark & \xmark & \cmark & \xmark \\
Streaming video LLMs & \xmark & \cmark & \xmark & \cmark & \cmark \\
BMVM (ours) & \cmark & \cmark & \cmark & \cmark & \cmark \\
\bottomrule
\end{tabular}

\end{table}

\paragraph{Cost model.}
Each timestep incurs cheap scan cost for lightweight features and query cost for VLM calls. The default benchmark fixes the number of VLM calls per timestep; secondary sweeps fix total GPU-seconds per hour. This separation lets the benchmark compare scheduling policy quality without conflating it with the raw speed of a particular VLM backend.
